\chapter{Свежая литература о суперсвязностных дефектах и проверка новизны}

Предыдущий гейт оставил одну внутреннюю лазейку: гладкое продолжение пары
$L/L^*$ может задаваться не чётной $SO(3)$-связностью, а нечётной
компонентой суперсвязности, теряющей обратимость на ядре. Перед новым
расчётом необходимо проверить первичную литературу и отделить уже известный
общий механизм от специфического содержания проекта.

Настоящая глава является литературно-архитектурным гейтом. Она не объявляет
сходство терминов доказательством и не использует свежие препринты как
подтверждение физических выводов.

\section{Иерархия источников}

Источники разделены по статусу.

\subsection{Фундаментальный проверенный результат}

Работа Канно и Сугимото ``Anomaly and Superconnection''
(\path{arXiv:2106.01591}, PTEP 2022) рассматривает фермионы с
пространственно-зависимой матричной массой. Две градуированные ветви несут
связности $A_+$ и $A_-$, а нечётное поле $T$ входит как внедиагональная
часть
\begin{equation}
 \mathcal A=
 \begin{pmatrix}
  A_+&iT^\dagger\\
  iT&A_-
 \end{pmatrix}.
 \label{eq:v5-literature-standard-superconnection}
\end{equation}
Чернов характер этой суперсвязности выражает аномалию, а асимптотические
данные матричной массы связаны с индексами типа Каллиаса.

Для нашей ветви это устанавливает допустимый язык:
\begin{equation}
 L/L^*\leftrightarrow E_+/E_- ,\qquad
 \text{переход}\leftrightarrow T,\qquad
 \text{дефект}\leftrightarrow\operatorname{rank}T\text{ падает}.
\end{equation}
Но работа не выводит конкретный $T$ из $M_{35}$ и не фиксирует его
энергетический функционал.

\subsection{Свежая конкретная реализация 2026 года}

В препринте Raymond--Järvinen--Kiritsis--Nitti--Préau
``The Tachyon Chern--Simons action with a generic tachyon field, and
baryons in V-QCD'' (\path{arXiv:2607.24915}, 27 июля 2026 года) поле $T$
является произвольной комплексной матрицей между левой и правой
калибровочными ветвями. Дефект определяется как множество
\begin{equation}
 \mathcal Z=\{x:\det T(x)=0\}.
 \label{eq:v5-literature-noninvertibility-locus}
\end{equation}
Чернов характер суперсвязности локализуется около множества потери
обратимости, а топологический заряд выражается классом Черна
суперсвязности.

Это наиболее важная свежая подсказка для проекта. Гладкое ядро не обязано
пониматься как продолжение отдельной линии. Вне ядра нечётный оператор
обратим и определяет разложение на спектральные ветви; на ядре его ранг
падает, поэтому сами ветви перестают быть отдельными расслоениями.

Однако в указанной работе динамика не возникает из одного топологического
характера. Модель заранее содержит пятимерные калибровочные поля,
DBI-сектор, голографический фон и тахионное действие. Суперсвязность
жёстко организует топологический и аномальный сектор, но не заменяет весь
энергетический функционал.

\subsection{Свежий операторно-алгебраический язык}

Работа Jia--Luo--Tian--Wang--Zhang ``Categorical Symmetries via Operator
Algebras'' (\path{arXiv:2604.25821}, 28 апреля 2026 года) описывает
скрученные категориальные симметрии через $C^*$-алгебры расслоений Фелла
над группоидами. Волокна над морфизмами являются соответствиями, а
скручивание кодируется центральным $U(1)$-расширением и групоидным
когомологическим классом.

Более специальная работа Bédos--Conti ``Actions of Fell bundles''
(\path{arXiv:2509.10822}) строит $C^*$-соответствия из действий расслоений
Фелла.

Эти результаты подтверждают корректность языка
\begin{equation}
 \text{стрелка}\longmapsto\text{соответствие с линейным скручиванием},
\end{equation}
но не дают локального поля $T(x)$, конечной энергии или физического
спектра. Их роль --- топологическая типизация переходов и аномалий.

\subsection{Близкий, но непроверенный сосед}

Препринт P.~Ogonowski ``Self-Reconstructing Codazzi Defects,
$\mathbb{CP}^1$ Quantization, and the Minimal Standard-Model Carrier''
(\path{arXiv:2606.12693v2}) использует разрешённую ссылку
$\mathbb{CP}^1$, линию степени один и реконструкцию конечного носителя.
Тематическое пересечение с проектом очень велико.

Но его статус необходимо записать строго:
\begin{enumerate}
 \item это свежий непрошедший рецензирование препринт в категории
 \texttt{physics.gen-ph};
 \item версия arXiv v2 датирована 25 июня 2026 года;
 \item завершение Дирака--Каллиаса, Рисса и Schur--Berry сформулировано
 как последующий этап, а не как уже выполненный результат;
 \item сходство архитектуры не является независимым экспериментальным или
 математическим подтверждением нашей ветви.
\end{enumerate}

Тем самым исправляется предварительное слишком сильное чтение: этот
препринт показывает актуальность направления, но не закрывает наш
операторный и энергетический разрыв.

\section{Что в нашей конструкции уже известно литературе}

По отдельности не являются новыми:
\begin{enumerate}
 \item градуированный пакет с двумя связностями и нечётным матричным полем;
 \item дефект как множество нулей или потери обратимости нечётного поля;
 \item индекс типа Каллиаса из асимптотической матричной массы;
 \item чернов характер суперсвязности как носитель топологического заряда;
 \item соответствия и линейные скручивания Фелла над группоидом.
\end{enumerate}

Следовательно, дальнейший текст не должен заявлять изобретение
``суперсвязностной частицы'' вообще.

\section{Где может находиться специфическая новизна проекта}

Потенциально специфична не отдельная техника, а их жёсткая сборка:
\begin{enumerate}
 \item ориентационная степень выводится из неравноранговых углов
 $p_{20}-p_{15}$;
 \item переходам $E/E^*$ функториально назначаются $L/L^*$ без нового
 фиксированного семейного дублета;
 \item эта конструкция обязана быть совместима с уже замороженными
 $H_{15}$, KO6, физическим чтением и единым следом $M_{35}$;
 \item все неудачные способы получить энергию, индекс и нормировку
 сохраняются в одном проверяемом реестре.
\end{enumerate}

Эта новизна пока является только кандидатом. Она станет содержательной
лишь после построения правильно типизированного нечётного оператора и
проверки его индекса без добавления нового носителя.

\section{Новый центральный тест рангов}

Литературная переформулировка немедленно выявляет различие между
стандартным механизмом и нашим родителем. В
\eqref{eq:v5-literature-standard-superconnection} поле $T$ обычно является
квадратной матрицей
\begin{equation}
 T:E_+\longrightarrow E_- ,\qquad
 \rank E_+=\rank E_-.
\end{equation}
Тогда оно может быть обратимым вне компактного дефекта и терять ранг
только на ядре.

В связывающей алгебре $M_{35}$ естественная нечётная компонента имеет вид
\begin{equation}
 T\in E=M_{20\times15}(\mathbb C),
 \qquad T:\mathbb C^{15}\longrightarrow\mathbb C^{20}.
 \label{eq:v5-literature-rectangular-odd-field}
\end{equation}
Она нигде не является обратимой. Даже при максимальном ранге остаётся
пятимерный коядровый сектор:
\begin{equation}
 \dim\operatorname{coker}T\geq20-15=5.
\end{equation}
Поэтому множество ``$det T=0$'' для исходного угла не определено, а
потеря обратимости не локализует ядро: необратимость присутствует всюду.

\begin{nogocriterion}[Прямоугольный нечётный угол]
Нельзя напрямую перенести квадратный тахионный или матрично-массовый
механизм на угол $20\times15$. Пока не выведен равногранный пакет,
нечётная компонента $M_{35}$ не может быть обратима вне дефекта и не задаёт
локализованное множество потери обратимости.
\end{nogocriterion}

Существуют две формальные возможности:
\begin{enumerate}
 \item ограничить семейный угол на физический ранг пятнадцать и получить
 квадрат $15\times15$;
 \item работать на паре пространства стрелок и сопряжённых стрелок
 $E\oplus E^*$, то есть на квадратном носителе ранга $300+300$.
\end{enumerate}
Ни одна из них пока не выведена как допустимая нечётная компонента
существующего родителя. Первая может выбросить канонический опорный сектор,
вторая может скрыто удвоить носитель и след.

\section{Изменение исследовательского маршрута}

Старая постановка требовала сначала построить гладкую чётную
$SO(3)$-связность, компенсирующую энергию ежа. Свежая литература предлагает
более первичный вопрос:
\begin{equation}
 \boxed{
 \begin{gathered}
 \text{существует ли нечётный }T(x),\\
 \text{обратимый вне ядра и теряющий ранг только на ядре?}
 \end{gathered}}
\end{equation}

Если такой оператор выводится из текущей градуированной архитектуры, то:
\begin{enumerate}
 \item $L/L^*$ читаются как его асимптотические спектральные линии;
 \item центр дефекта является не местом продолжения каждой линии, а
 местом разрушения спектрального разложения;
 \item индекс вычисляется из асимптотического класса суперсвязности;
 \item вопрос энергии ставится затем для того же $T$, а не ремонтируется
 отдельным монополем.
\end{enumerate}

Но топология не фиксирует энергию автоматически. Даже самые свежие
реализации используют отдельные DBI, Янга--Миллса, голографические или
потенциальные данные. Поэтому успешный нечётный индексный оператор будет
лишь следующим математическим мостом, а не физическим замыканием.

\section{Вердикт и следующий гейт}

Литературный аудит не обнаружил готовой теории, которую можно просто
перенести в проект. Он обнаружил более точный научный язык и одновременно
новое ранговое препятствие.

\begin{equation}
 \boxed{
 \begin{gathered}
 \text{язык нечётной суперсвязности и потери обратимости: стандартен},\\
 \text{хопфово-моритова сборка проекта: потенциально специфична},\\
 \text{прямой угол }20\times15:\text{ не локализует дефект},\\
 \text{необходим равногранный оператор без скрытого расширения}.
 \end{gathered}}
\end{equation}

Следующий гейт должен классифицировать обе равногранные возможности и
проверить, существует ли на них KO6-совместимый нечётный оператор с
асимптотическими линиями $L/L^*$ и единичным индексом:
\path{version5_hopf_pair_odd_core_extension_gate.tex}.

\begin{protocol}
\begin{enumerate}
 \item первичные источники проверены: \textbf{да};
 \item рецензированные и свежие препринты разделены: \textbf{да};
 \item Codazzi-препринт считается завершением нашей ветви: \textbf{нет};
 \item нечётная суперсвязность как общий язык нова: \textbf{нет};
 \item дефект как потеря обратимости нов: \textbf{нет};
 \item ориентационный функтор $p_{20}-p_{15}\mapsto L/L^*$:
 \textbf{кандидат специфического результата};
 \item прямоугольный угол $20\times15$ обратим вне ядра: \textbf{нет};
 \item постоянный минимальный коядровый ранг: \textbf{пять};
 \item равногранный оператор выведен: \textbf{нет};
 \item энергия следует из топологического характера: \textbf{нет};
 \item физическое замыкание: \textbf{не достигнуто}.
\end{enumerate}
\end{protocol}

Структурированный сертификат сохранён в
\path{s2t_v5_recent_superconnection_defect_literature_novelty_gate_results.json}.